\subsectionaddtolist{الف) \lr{External Network}}

وقتی با \lr{docker compose} یک پروژه را بالا می‌آوریم، داکر به‌صورت پیش‌فرض برای همان پروژه یک شبکه‌ی اختصاصی می‌سازد و با پایین آوردن پروژه آن را حذف می‌کند. شبکه‌ی \lr{external} یعنی شبکه‌ای که بیرون از این چرخه‌ی عمر ساخته شده است؛ ما فقط اعلام می‌کنیم که از قبل وجود دارد و داکر باید سرویس را به آن وصل کند، نه اینکه آن را بسازد یا پاک کند.

\begin{lstlisting}[language=Yaml]
# docker network create network-name
networks:
  network-name:
    external: true
\end{lstlisting}

کاربرد اصلی‌اش این است که چند پروژه‌ی \lr{compose} مستقل بتوانند همدیگر را با نام سرویس صدا بزنند، بدون اینکه مجبور باشیم همه را در یک فایل واحد جمع کنیم. نکته‌ی مهم این است که چون داکر کامپوز فعلی مالک آن شبکه نیست، اگر شبکه از قبل ساخته نشده باشد کامپوز خطا می‌دهد و پاک کردن آن هم وظیفه‌ی خود ماست.

\subsectionaddtolist{ب) \lr{Linting}}

لینتینگ یعنی تحلیل ایستای کد؛ یعنی بررسی کد بدون اجرای آن برای پیدا کردن خطاهای احتمالی، الگوهای مشکوک و ناسازگاری با قواعد سبک تیم. تفاوتش با کامپایلر این است که کامپایلر فقط می‌پرسد آیا این کد معتبر است؟ ولی لینتر می‌پرسد آیا این کد خوب است؟: متغیر استفاده‌نشده، خطای نادیده‌گرفته‌شده، ایمپورت اضافی، نام‌گذاری خلاف قرارداد، پیچیدگی بیش از حد یک تابع و مواردی از این دست.

معمولا دو لایه دارد: \lr{formatter} که فقط ظاهر کد را یکدست می‌کند (مثل \lr{gofmt} یا \lr{prettier}) و \lr{linter} که به معنا و ساختار کد کار دارد (مثل \lr{golangci-lint} یا \lr{ESLint}). ارزش اصلی‌اش در تیم این است که بحث‌های سلیقه‌ای از \lr{code review} حذف می‌شود و بازبین به‌جای شمردن فاصله‌ها، روی منطق تمرکز می‌کند. به همین دلیل لینت تقریبا همیشه اولین و ارزان‌ترین مرحله‌ی پایپلاین است.

البته بعضا مرحله فرمت و لینت را از هم جدا می‌کنند، اما بعضی تیم‌ها هم هر دو را در یک مرحله انجام می‌دهند.


\subsectionaddtolist{ج) \lr{Image Registry}}

رجیستری، انبار نسخه‌‌بندی‌شده‌ی \lr{image}ها است؛ همان نقشی که \lr{npm} برای پکیج‌های جاوااسکریپت یا \lr{Maven Central} برای جاوا دارد. ایمیج با نام و تگ در آن ذخیره می‌شود، با \lr{docker push} بالا می‌رود و با \lr{docker pull} پایین می‌آید.

نکته‌ی فنی مهم این است که رجیستری ایمیج را به‌صورت یک فایل یکپارچه نگه نمی‌دارد، بلکه یک \lr{manifest} به‌علاوه‌ی لایه‌های ایمیج را ذخیره می‌کند و هر لایه با \lr{digest} شناخته می‌شود؛ بنابراین لایه‌های مشترک بین ایمیج‌ها فقط یک بار ذخیره و منتقل می‌شوند. همین \lr{digest} است که اجازه می‌دهد به‌جای تگ به \lr{image@sha256:...} ارجاع دهیم.
